% Resume: Arumalla Mohit Krishna
% Compile with pdfLaTeX.

\documentclass[letterpaper,10pt]{article}

\usepackage[margin=0.5in]{geometry}
\usepackage[T1]{fontenc}
\usepackage{lmodern}
\usepackage{microtype}
\usepackage{titlesec}
\usepackage{enumitem}
\usepackage[hidelinks]{hyperref}
\usepackage{tabularx}
\input{glyphtounicode}

\pdfgentounicode=1
\pagestyle{empty}
\urlstyle{same}
\raggedright
\setlength{\parindent}{0pt}
\setlength{\tabcolsep}{0pt}

\titleformat{\section}
  {\large\scshape}{}{0em}{}[\titlerule]

\titlespacing*{\section}{0pt}{6pt}{3pt}

\newcommand{\entry}[2]{%
  \noindent
  \begin{tabularx}{\linewidth}
    {@{}>{\raggedright\arraybackslash}X@{\hspace{6pt}}r@{}}
    \textbf{#1} & #2
  \end{tabularx}\par
}

\newcommand{\bullets}{
  \begin{itemize}[
    leftmargin=1.2em,
    topsep=2pt,
    itemsep=1pt,
    parsep=0pt,
    partopsep=0pt
  ]
}

\newcommand{\finishbullets}{
  \end{itemize}
}

\begin{document}

%----------HEADING----------
\begin{center}
    {\LARGE\bfseries Arumalla Mohit Krishna}\\[3pt]
    +1 (631) 352-8643 $|$
    \href{mailto:mohitkrishna.arumalla@stonybrook.edu}
    {mohitkrishna.arumalla@stonybrook.edu} $|$
    \href{https://linkedin.com/in/mohit-krishna-arumalla}
    {LinkedIn}
\end{center}
\vspace{-8pt}

%-----------EDUCATION-----------
\section{Education}

\entry
  {\textbf{Stony Brook University}}
  {Aug 2025 -- Dec 2026}

M.S. in Computer Science \hfill New York, USA\\
\textit{
Coursework: Natural Language Processing,
Introduction to Computer Vision, Vision Language Modeling
}
\vspace{3pt}

\entry
  {\textbf{Birla Institute of Technology and Science, Pilani}}
  {Aug 2019 -- 2023}

B.E. in Electronics \& Communication; GPA: 8.02/10 \hfill India
%-----------PROFESSIONAL SUMMARY-----------
\section{Professional Summary}

Software engineer with 2+ years of experience engineering \textbf{distributed systems, real-time data pipelines}, and \textbf{semantic search} capabilities, currently pursuing an M.S. in Computer Science. Experienced in building \textbf{robust data infrastructures} and scalable retrieval systems, with hands-on expertise in \textbf{agent-oriented design patterns}, GPU model serving, and \textbf{end-to-end data workflows}. Proficient in Python, Go, PyTorch, and cloud-native infrastructure, with a strong focus on \textbf{systematic reasoning}, \textbf{ownership}, and complex problem decomposition.

%-----------TECHNICAL SKILLS-----------
\section{Technical Skills}

\textbf{Languages:}
Python, Go, C++, SQL, Java\\

\textbf{Machine Learning \& Agents:}
PyTorch, TensorFlow/Keras, Transformers, Self-Supervised Learning, LoRA, Text Embeddings, Vector Search, Agent Context Engineering, Long-Horizon State Management\\

\textbf{Data \& Systems:}
gRPC, Kafka, Apache Flink, Flink SQL, Debezium/CDC,
PostgreSQL, Redis, Elasticsearch, Airflow, BigQuery\\

\textbf{Infrastructure:}
NVIDIA Triton, AWS, Docker, Kubernetes, Linux, Git,
CI/CD, Istio/Envoy, New Relic

%-----------EXPERIENCE-----------
\section{Experience}

\entry
  {\textbf{Software Engineer II} $|$ \textbf{Infra.Market, Bangalore}}
  {Apr 2024 -- Aug 2025}

\bullets

  \item Developed \textbf{semantic search} capabilities with Amazon Titan text embeddings
  and Elasticsearch vector search, enabling sophisticated natural-language product
  discovery through conversational chat interfaces.

  \item Engineered and maintained a \textbf{CDC pipeline} combining Flink, Kafka, Debezium,
  PostgreSQL, and Elasticsearch; reducing synchronization latency from 10--15 minutes to approximately 500 ms with \textbf{exactly-once} processing semantics, critical for high-throughput environments.

  \item Deployed and optimized \textbf{service mesh} across 15+ microservices with Istio/Envoy, integrating distributed tracing via New Relic to enhance monitoring, reliability, and observability in production.

\finishbullets

\entry
  {\textbf{Software Engineer I} $|$ \textbf{Infra.Market, Chennai}}
  {Jul 2023 -- Mar 2024}

\bullets

  \item Owned highly-available Search and Location microservices; designed and deployed \textbf{gRPC APIs} processing 500K+ requests daily for product discovery and geospatial lookups.

  \item Addressed complex PostgreSQL connection leaks within Gorm ORM, reducing connection-pool exhaustion by 85\%; validated stability using load testing at 10K QPS.

  \item Built and optimized Airflow pipelines for BigQuery-to-Elasticsearch data migration with \textbf{dynamic filtering} and automated monitoring, ensuring low-latency updates to support real-time search index availability.

\finishbullets

%-----------PROJECTS-----------
\section{Projects}

\entry
  {GPU-Optimized Personalized Product Search}
  {Aug 2026 -- Oct 2026}

\textit{Python, PyTorch, BGE, NVIDIA Triton, Go, gRPC}

\bullets

  \item Developed two-tower retrieval architecture using \texttt{BAAI/bge-small-en-v1.5} embeddings
  for query and product spaces, combined with a trainable user tower to incorporate interaction history for personalized search.

  \item Designed and implemented an efficient inference pipeline featuring precomputed product embeddings,
  approximate nearest neighbor retrieval, and a Go/gRPC gateway integrated with NVIDIA Triton for scalable GPU-backed query encoding.

  \item Benchmarked batching and concurrency scenarios under steady and bursty traffic, monitoring latency, throughput,
  GPU utilization, and costs to optimize for production readiness; evaluated Recall@K and NDCG@K to ensure retrieval quality.

\finishbullets

\entry
  {Attention-Based Hallucination Detection in NMT}
  {Apr 2026 -- May 2026}

\textit{Python, PyTorch, mBART, Transformers, scikit-learn}

\bullets

  \item Built an English-German hallucination detection pipeline based on teacher-forced mBART attention features using WMT22/WMT23 quality estimation data.

  \item Extracted cross- and self-attention features from 12 decoder layers and 16 heads, considering metrics like entropy, mass concentration, variance,
  and grounding ratios to capture hallucination signals.

  \item Trained classifiers at token, sentence, and head levels achieving \textbf{0.712 F1} and \textbf{0.695 ROC-AUC} on a 50K-example dataset;
  proposed hallucination-aware reranking leveraging attention-based localization and risk scoring.

\finishbullets

\entry
  {Prompt/LoRA for Self-Supervised Learning}
  {Ongoing}

\textit{Python, PyTorch, Vision Transformers}

\bullets

  \item Investigating domain-specific visual prompting and LoRA parameter-efficient adaptation in a self-supervised ViT framework to learn transferable representations across diverse image domains.

  \item Evaluating backbone performance without prompts or LoRA at inference time; designing comparative analyses for efficient parameter use and domain generalization.

\finishbullets

\end{document}